\scanpage{7}
\deleted{Pour tout espace analytique $S$, soit $F(S)$ l'ensemble \inserted{\unsure{quotient} des objets suivants} des \uline{foncteurs} qui, à chaque $V\in\operatorname{Rep}(G)$, associe}

Soit $S$ un espace analytique. Une $G$-\uline{subordination} \inserted{\uline{rigidifiée}} \uline{sur} $S$ est la donnée d'un $\otimes$-foncteur qui, à tout $V\in\operatorname{Ob}\operatorname{Rep}(G)$, associe une structure de famille de Hodge \inserted{mod isogénie} sur $V_{\mathbb{C}}\times S\unsure{\underset{\mathrm{D}}{=}|\underset{\delta}{=}|\underset{\mathrm{Df}}{=}}V_S$ \deleted{polarisables en chaque point,} [i.e. une bigraduation \inserted{analytique \uline{réelle}} sur $V_{\mathbb{C}}\times S$, \inserted{\deleted{\unsure{autres}} avec $V_S^{pq}=\overline{V_S^{q,p}}$,} telle que la $p$-filtration correspondante soit une filtration analytique complexe] et telle que \inserted{a) degré total soit celui défini par $i$ b) $\mathbb{Q}(-1)$ devienne \unsure{St|Sf} \unsure{de H.} de bidegré $(1,1)$, et} \inserted{c)} pour \inserted{$\forall$ $V$ homogène de degré $n$ et} $\forall\,s\in S$, la fibre en $s$ soit polarisable \deleted{\unsure{pour} $\forall$ \unsure{p\textsuperscript{t} fixé}}\inserted{à l'aide d'une forme $\varphi:V\otimes V\to\mathbb{Q}(-n)$ \inserted{qui soit $G$-invariante (NB. si $S$ est connexe, il suffit de l'exiger cette condition en \uline{un} point \deleted{\unsure{$s$}} de $S$, car la forme restera de polarisation partout, puisqu'elle est non dég. partout}}

\deleted{\unsure{Retenant|Retenons}}\inserted{Oubliant pour le moment la condition de} holomorphie de la filtration, on trouve :

a) pour $\forall$ p\textsuperscript{t} $s\in S$, on veut une façon \deleted{de faire} fonctorielle de mettre une structure de Hodge \inserted{(mod isogénie)} polarisable sur les $V\in\operatorname{Ob}\operatorname{Rep}(G)$, i.e. donnée par un $j=J(s):\unsure{\mathbb{S}|\mathrm{S}}\longrightarrow G_{\mathbb{R}}$ satisfaisant les conditions du th, i.e. $j'=J'(s):\mathbb{U}\unsure{\rightharpoonup|\to}G'_{\mathbb{R}}$ satisfaisant $j'(\deleted{\unsure{C}}-1)=\eta$, et tel que $\operatorname{Zent}_{G_{\mathbb{R}}^{\prime\,0}}(j'(s)(C))$ soit compact maximal.

b) Pour $s$ variable, $j'(s)$ dépend de $s$ de façon analytique réelle.
% archive sheet number: 6

